\documentclass[10pt,a4paper]{amsart}
\usepackage[margin=19mm]{geometry}
\usepackage{amsmath,amssymb,mathtools}
\usepackage[scr=rsfs,cal=euler]{mathalfa}
\usepackage{xfrac}
\usepackage[unicode,hidelinks,bookmarks=false]{hyperref}
\newcommand{\ie}{i.e.\ }
\newcommand{\cf}{cf.\ }
\newcommand{\resp}{respectively\ }
\newcommand{\Subsection}{\subsection}
\providecommand{\nobreakdash}{\nobreak\hbox{-}}
\setlength{\emergencystretch}{2em}
\multlinegap0pt
\usepackage{tikz}
\usepackage{amssymb}
\usepackage{bm}
\usepackage[shortlabels]{enumitem}
\newtheorem{lemma}{Lemma}
\newtheorem{theorem}{Theorem}
\newtheorem{claim}{Claim}
\newcommand{\E}{\mathbb E}
\newcommand{\R}{\mathbb R}
\newcommand{\ceil}[1]{\left\lceil #1 \right\rceil}
\newcommand{\fS}{\mathcal S}
\newcommand{\fX}{\mathcal X}
\newcommand{\paren}[1]{\left( #1 \right)}
\title{Separators for intersection graphs of spheres}
\author{Jacob Fox, Jonathan Tidor}
\date{Source: arXiv:2603.22204v1 (CC BY 4.0)}
\begin{document}
\maketitle

\noindent\emph{Source and licence.} Separators for intersection graphs of spheres. Version arXiv:2603.22204v1 (CC BY 4.0), distributed by arXiv under the Creative Commons Attribution 4.0 International licence, http://creativecommons.org/licenses/by/4.0/, as recorded in the arXiv licence metadata for this exact version and shown on the abstract page https://arxiv.org/abs/2603.22204v1. Full licence notice: \texttt{licenses/arxiv-2603.22204-CC-BY-4.0.txt}.\par\smallskip

\noindent\textit{Editorial scope.} Counted unit: the article's theorem thm:sphere-separator-main (source0.tex lines 207--211). The article's complete original proof of it is delivered with it (source0.tex lines 226--337, one \texttt{proof} environment of 112 lines, which the article places away from the statement). No mathematics is written, repaired or re-derived: the statement and the proof are the article's own lines, in the article's own notation, and no retained line is substituted or re-flowed.  Retained uncounted as the source-local context the proof needs: lines 132--135; lines 215--218.  Nothing was omitted from any retained span, so the package carries no omission record.  No retained line contains a citation, so the wrapper carries no bibliography and no bibliography entry.  The retained lines contain the article's own diagram code, which the wrapper typesets with the standard tikz packages; no image file is delivered and the package declares no asset dependency.  Editorial formatting only, in addition: an AMS \texttt{amsart} preamble that restates the article's own declarations and macros used by the retained lines, namely the environments \texttt{lemma}, \texttt{theorem}, \texttt{claim} on separate counters, together with the standard packages those lines need.  The retained environments print in this document as Lemma 1, Theorem 1, Claim 1 in order of appearance, whereas the article prints the same items with the numbering of its own full document.  The complete native source of the article as distributed in the arXiv e-print for this exact version is bundled unchanged as \texttt{sources/197005/source0.tex}.
\begin{lemma}
\label{lem:balancing-separators}
Suppose graph $G=(V,E)$ is such that every induced subgraph has a $c$-balanced separator of size at most $s$. Then $G$ has a balanced separator of size at most $\lceil \log_c(2/3)\rceil s$.
\end{lemma}

\begin{lemma}
\label{lem:nested-sphere-separator}
Let $\fS$ be a collection of spheres in $\R^d$. Suppose that $G(\fS)$ has maximum degree at most $\Delta$ and that there exists a point $p\in\R^d$ contained in the interior of at least $2s$ spheres. Then there exists a set $\fX\subseteq\fS$ with $|\fX|\leq\Delta$ such that every connected component of $G(\fS)\setminus\fX$ has size at most $|\fS|-s$.
\end{lemma}

\begin{theorem}
\label{thm:sphere-separator-main}
Let $\fS$ be a collection of $n$ spheres in $\R^d$. Then $G(\fS)$ has a balanced separator of size at most
\[O_d\paren{\paren{\sum_{S\in\fS}\deg(S)^{\tfrac1{d-1}}}^{1-\tfrac1d}}.\]
\end{theorem}

\begin{proof}[Proof of Theorem~\ref{thm:sphere-separator-main}]
Set \[\Sigma=C_d\paren{\sum_{S\in\fS} \deg(S)^{\tfrac1{d-1}}}^{1-\tfrac1d}.\] Choosing $c_d>0$ small in terms of $d$ and $C_d$ large in terms of $c_d$, we will construct a $(1-c_d)$-balanced separator for $G(\fS)$ of size at most $\Sigma$. By Lemma~\ref{lem:balancing-separators}, losing another constant factor depending only on $d$, this suffices to prove the desired result. Assume $\Sigma<n$, since the result is trivial otherwise.

Through the proof, if we have not already found the desired separator, we will define sets $\fX_i \subset \fS$ for $i=1,2,3,4$ with $|\fX_i| \leq \Sigma/4$ for each $i$ so that $\fX_1\cup \fX_2\cup \fX_3\cup \fX_4$ is the desired separator.

\vspace{5pt} \noindent \textbf{Step 1:}
First, define $\fX_1\subset \fS$ to be the $\Sigma/4$ vertices of highest degree in $G(\fS)$. Let $\Delta$ be the maximum degree of the resulting graph $G(\fS)\setminus\fX_1$. We claim that $\Delta\leq\Sigma/4$ and $\Delta\leq c_dn$. If the first does not hold, then $G(\fS)$ would have at least $\Sigma/4$ vertices of degree at least $\Sigma/4$, meaning that
\[\Sigma=C_d\paren{\sum_{S\in\fS} \deg(S)^{1/(d-1)}}^{1-1/d}\geq C_d\paren{(\Sigma/4)(\Sigma/4)^{1/(d-1)}}^{1-1/d}=\frac{C_d}4\Sigma.\]
Clearly this is impossible for $C_d>4$. Similarly, if the second fails then $G(\fX)$ would have at least $\Sigma/4$ vertices of degree at least $c_dn$, meaning that
\[\Sigma\geq C_d\paren{(\Sigma/4)(c_dn)^{1/(d-1)}}^{1-1/d}=\frac{C_dc_d^{1/d}}{4^{1-1/d}}\Sigma^{1-1/d}n^{1/d}>\frac{C_dc_d^{1/d}}{4^{1-1/d}}\Sigma,\]
where the final inequality holds since $\Sigma<n$. Again, this is impossible for $C_d$ sufficiently large in terms of $c_d$.

\vspace{5pt} \noindent \textbf{Step 2:} Write $\fS_1=\fS\setminus \fX_1$. The goal of this step is to pass to a subset $\fS_2\subseteq \fS_1$ such that every sphere $S\in\fS_2$ contains fewer than $4c_dn$ centers of spheres in $\fS_2$. Let $S_0\in\fS_1$ be a sphere of minimal radius that contains $4c_dn$ centers of spheres in $\fS_1$. (If no such sphere exists, taking $\fS_2=\fS_1$ completes this step.) Define $\fS_{\supseteq S_0}\subseteq \fS_1$ to be the set of spheres which contain $S_0$ entirely within their interior. First suppose that $|\fS_{\supseteq S_0}|\geq 2c_dn$. In this case we can apply Lemma~\ref{lem:nested-sphere-separator} to find a separator $\fX_2\subset\fS_1$ such that $|\fX_2|\leq\Delta\leq\Sigma/4$ and each connected component of $G(\fS_1)\setminus\fX_2$ has size at most $(1-c_d)n$. Then $\fX_1\cup\fX_2$ is a $(1-c_d)$-balanced separator for $G(\fS)$ of size at most $\Sigma/2$.

Thus we can assume that $|\fS_{\supseteq S_0}|< 2c_dn$. In this case, define $\fX_2=N(S_0)$, the neighborhood of $S_0$ in $G(\fS)$. Define $\fS_2\subseteq\fS_1$ to be the set of spheres which are entirely contained within $S_0$. Clearly $|\fX_2|\leq \Delta\leq\Sigma/4$.

\begin{claim}
\label{claim:sphere-separator-proof}
If $G(\fS)\setminus(\fX_1\cup \fX_2)$ has a connected component of size more than $(1-c_d)n$, then it is also a connected component of $G(\fS_2)$. 
\end{claim}
\begin{proof}
Every connected component of $G(\fS)\setminus (\fX_1\cup \fX_2)$ is either made up of spheres entirely inside $S_0$ or entirely outside $S_0$. The former are also connected components of $G(\fS_2)$ while the latter are not. There are at least $4c_dn$ centers of spheres of $\fS_1$ inside of $S_0$. Since at most $2c_dn$ of these correspond to spheres surrounding $S_0$ and at most $\Delta$ correspond to spheres which intersect $S_0$, we see that there are at least $2c_dn-\Delta\geq c_dn$ spheres inside $S_0$. Therefore any connected component outside of $S_0$ has size at most $(1-c_d)n$.
\end{proof}

We now construct a separator for $G(\fS_2)$. We know that every sphere in $\fS_2$ contains fewer than $4c_dn$ centers of spheres in $\fS_2$. We can also assume that $|\fS_2|\geq(1-c_d)n$, otherwise $\fX_1\cup\fX_2$ is already the desired $(1-c_d)$-balanced separator.

\vspace{5pt} \noindent \textbf{Step 3:} Let $B_{\mathrm{in}}$ be a ball of minimal radius which contains at least $4c_dn$ centers of spheres in $\fS_2$. Since we showed that no sphere in $\fS_2$ contains $4c_dn$ centers, we see that no sphere in $\fS_2$ contains $B_{\mathrm{in}}$ entirely in its interior. This means that each of the $4c_dn$ spheres whose centers lie in $B_{\mathrm{in}}$ also intersect $B_{\mathrm{in}}$.

Write $R$ for the radius of $B_{\mathrm{in}}$ and define $B_{\mathrm{out}}$ to be the ball with the same center and radius $2R$. By the minimality of $B_{\mathrm{in}}$, we see that $B_{\mathrm{out}}$ contains at most $4^d\ceil{4c_dn}$ centers. In particular, choosing $c_d$ sufficiently small, there are at least $c_dn$ sphere centers outside of $B_{\mathrm{out}}$. Let $S_{\mathrm{rand}}$ be a sphere with the same center as $B_{\mathrm{in}}$ and $B_{\mathrm{out}}$ whose radius is chosen uniformly at random in $[R,2R]$.

For a sphere $S\in\fS$, write $r(S)$ for the radius of $S$. Write $\tilde r(S)$ for the radius of the spherical cap formed by intersecting the ball bounded by $S$ with $S_{\mathrm{rand}}$. (The radius of a spherical cap is the minimal $\tilde r(S)$ so that the cap is contained in a ball of radius $\tilde r(S)$.)

Choose some $C_d'$ sufficiently large in terms of $d$ and let $\fX_3$ consist of the spheres $S\in \fS_2$ which intersect $S_{\mathrm{rand}}$ and which satisfy $\tilde r(S)\leq C_d' R(\deg(S)/\Sigma)^{1/(d-1)}$.

\begin{claim}
\label{claim:expected-x3-size}
$\E[|\fX_3|]\leq\Sigma/4$.
\end{claim}
\begin{proof}
We need to upper bound the probability that a sphere $S$ intersects $S_{\mathrm{rand}}$ in a spherical cap of radius at most $\tilde r$. If $\tilde r\leq r(S)$, the probability is upper bounded by 
\[\frac{2\paren{r(S)-\sqrt{r(S)^2-\tilde r^2}}}R\leq \frac{2\tilde r}{R}.\]
To see this, note that this event is the same as the event that $S_{\mathrm{rand}}$ intersects the line connecting the center of $B_{\mathrm{in}}$ and the center of $S$ in one of the bold segments in Figure~\ref{fig:expected-x3-size}.

Now if $\tilde r>r(S)$, the probability we need to upper bound is just the probability that $S_{\mathrm{rand}}$ and $S$ intersect. This quantity is at most $2r(S)/R<2\tilde r/R$. 

Since $\fX_3$ is defined with the cutoff $\tilde r = C_d' R(\deg(S)/\Sigma)^{1/(d-1)}$, for each $S\in\fS_2$ we have the bound
\[\Pr\left[S\text{ intersects }S_{\mathrm{rand}}\text{ and }\tilde r(S)\leq C_d' R(\deg(S)/\Sigma)^{1/(d-1)}\right]\leq 2C_d'(\deg(S)/\Sigma)^{1/(d-1)}.\]
This implies that
\begin{align*}
\E[|\fX_3|]
&\leq \sum_{S\in\fS_2} 2C_d'(\deg(S)/\Sigma)^{1/(d-1)}
=\frac{2C_d'}{\Sigma^{1/(d-1)}C_d^{d/(d-1)}}\Sigma^{d/(d-1)}
\leq \Sigma/4,
\end{align*}
where the equality is by the definition of $\Sigma$ and the last inequality holds for $C_d$ chosen appropriately large in terms of $C_d'$.
\end{proof}

Fix some choice of $S_{\mathrm{rand}}$ for which $|\fX_3|\leq\Sigma/4$. 

\begin{figure}
\centering
\begin{tikzpicture}[scale=4]
\clip (1.4-1,-0.5) rectangle + (2,1);

\draw(0,0) circle (1);
\draw(0,0) circle (2);

\draw(-2,0) -- (5,0);

\draw(1.4,0) circle (0.25);
\draw(1.4,0)--(1.4-0.6*0.25,0.8*0.25)--(1.4-0.6*0.25,0);
\draw(1.4,0)--(1.4+0.6*0.25,-0.8*0.25)--(1.4+0.6*0.25,0);
\draw[line width=2pt] (1.4-0.25,0)--(1.4-0.6*0.25,0);
\draw[line width=2pt] (1.4+0.25,0)--(1.4+0.6*0.25,0);


\node[fill=black, circle, inner sep=1pt] at (1.4,0) {};
\node[] at (1.4,-0.25) [below] {$S$};
\node[] at (1.4-0.6*0.25-0.03,0.8*0.25*0.5) [] {$\tilde r$};
\node[] at (1.4-0.6*0.25*0.5+0.09,0.8*0.25*0.5+0.03) [] {$r(S)$};

\end{tikzpicture}
\caption{\label{fig:expected-x3-size}
Bounding $\Pr[S\in\fX_3]$.}
\end{figure}

\vspace{5pt} \noindent \textbf{Step 4:} Define $\fS_3=\fS_2\setminus \fX_3$. We perform the following iterative process. Start with $\fX_4=\emptyset$ and $\fS_4=\fS_3$. We will define spheres $S_1,S_2,\ldots,S_\ell$. Suppose we have defined $S_1,\ldots, S_{i-1}$. Define $S_i$ to be a sphere in $\fS_4$ which intersects $S_{\mathrm{rand}}$ and has maximal radius among such spheres. Place the neighborhood $N(S_i)$ in $\fX_4$ and remove all spheres which intersect the ball bounded by $S_i$ from $\fS_4$ (in particular we remove $S_i$ and $N(S_i)$ from $\fS_4$). Repeat until no more spheres in $\fS_4$ intersect $S_{\mathrm{rand}}$.

We claim that the balls bounded by $S_1,S_2,\ldots,S_\ell$ are disjoint. For $i<j$ we know that $r(S_i)\geq r(S_j)$ and the ball bounded by $S_i$ is disjoint from the sphere $S_j$. The latter implies that $S_i,S_j$ do not intersect and that $S_i$ does not contain $S_j$; the former implies that $S_j$ does not contain $S_i$. 

\begin{claim}
\label{claim:x4-size}
$|\fX_4|\leq\Sigma/4$.
\end{claim}
\begin{proof}
By the definition of $\fX_3$, every sphere $S\in\fS_3$ which intersects $S_{\mathrm{rand}}$ also satisfies $\tilde r(S)>C_d'R(\deg(S)/\Sigma)^{1/(d-1)}$. Thus
\[|\fX_4|\leq\sum_{i=1}^\ell\deg(S_i)<\sum_{i=1}^\ell\Sigma\paren{\frac{\tilde r(S_i)}{C_d' R}}^{d-1}.\]
Now since the balls bounded by $S_1,\ldots,S_\ell$ are disjoint, we have
\[\sum_{i=1}^\ell \tau_{d-1}\tilde r(S_i)^{d-1}\leq \sigma_d (2R)^{d-1}\]
where $\tau_k,\sigma_k$ denote the $k$-volume of the unit ball in $\R^k$ and the $(k-1)$-volume of the unit sphere in $\R^k$, respectively. (In the above inequality, we use the fact that the $(d-1)$-volume of a spherical cap of radius $\tilde r$ on any sphere in $\R^d$ is at least the $(d-1)$-volume of a ball of radius $\tilde r$ in $\R^{d-1}$.) Thus we conclude that
\[|\fX_4|\leq \Sigma \frac{2^{d-1}\sigma_d}{C_d'^{d-1}\tau_{d-1}}\leq\Sigma/4,\]
where the last inequality holds for $C_d'$ chosen sufficiently large in terms of $d$.
\end{proof}

\begin{claim}
$\fX=\fX_1\cup \fX_2\cup \fX_3\cup \fX_4$ is a $(1-c_d)$-balanced separator for $G(\fS)$. 
\end{claim}
\begin{proof}
By Claim~\ref{claim:sphere-separator-proof}, it suffices to show that each connected component of $G(\fS_2)\setminus(\fX_3\cup \fX_4)$ has size at most $(1-c_d)n$. Note that any sphere in $\fS_2\setminus(\fX_3\cup \fX_4)$ which crosses $S_{\mathrm{rand}}$ must lie in the ball bounded by one of the $S_i$. Furthermore, no sphere in $\fS_2\setminus(\fX_3\cup \fX_4)$ crosses any of the $S_i$. Therefore any connected component of $G(\fS_2)\setminus(\fX_3\cup \fX_4)$ must either lie entirely within some sphere $S_i$, entirely within $S_{\mathrm{rand}}$, or entirely outside of $S_{\mathrm{rand}}$. We chose $\fS_2$ so that any sphere in $\fS_2$ contains at most $4c_dn$ centers of spheres in $\fS_2$. Thus the first type of connected component has size at most $4c_dn$. For the second type, note that $S_{\mathrm{rand}}$ is contained within $B_{\mathrm{out}}$ which we argued earlier contains at most $4^{d}\ceil{4c_dn}<(1-c_d)n$ centers of spheres in $\fS_2$. For the third type of connected component, we know that there are at least $4c_dn$ spheres which intersect $B_{\mathrm{in}}$. Since $B_{\mathrm{in}}$ is contained in the interior of $S_{\mathrm{rand}}$, we see that there are at most $(1-4c_d)n$ spheres outside of $S_{\mathrm{rand}}$, giving the desired bound on the third type of connected component.
\end{proof}
Since $|\fX|\leq\Sigma$, this completes the proof.
\end{proof}

\end{document}
